\chapter{Conclusie}\label{chap:conclusie}
Deze thesis demonstreert dat het mogelijk is om locatiegegevens on-device te obfusceren zonder de bruikbaarheid voor analyse significant aan te tasten. Door het ontwikkelen van een Android-library die locaties en tijdstempels vervormt via ruis en privacycirkels, wordt voorkomen dat individuele gebruikers geïdentificeerd kunnen worden via hun bewegingspatroon. De kern van de oplossing ligt in het behoud van data-utility: geaggregeerde inzichten, zoals mobiliteitstrends of drukteanalyses, blijven mogelijk, terwijl gevoelige individuele informatie wordt beschermd.

Een belangrijk succes is de volledige onafhankelijkheid van externe servers. Alle verwerking vindt plaats op het apparaat zelf, waardoor de originele, nauwkeurige locatie nooit het toestel verlaat. Dit elimineert het risico op datalekken bij overdracht of opslag. Technisch werd dit gerealiseerd door een modulair ontwerp, waarbij de library naadloos integreert met Android’s Fused Location Provider en WorkManager voor betrouwbare achtergrondverwerking. Daarnaast biedt de implementatie van een standaard obfuscatie-algoritme op basis van privacycirkels een evenwicht tussen anonymisering en nauwkeurigheid, geïnspireerd op eerder onderzoek.

De library verlaagt niet alleen de implementatiedrempel voor ontwikkelaars door complexe privacytechnieken te abstraheren, maar draagt ook bij aan een paradigmaverschuiving in databeheer. In plaats van privacy als een afterthought te behandelen, wordt bescherming een inherent onderdeel van het ontwerp. Dit sluit aan bij groeiende regelgevende eisen en maatschappelijke vraag naar transparante datapraktijken.

Hoewel verdere optimalisatie mogelijk blijft, toont deze aanpak aan dat technische oplossingen een sleutelrol kunnen spelen in het verzoenen van data-innovatie en individuele rechten. Door privacy-by-design toegankelijk te maken, zet deze thesis een stap richting een toekomst waarin technologie niet ten koste gaat van gebruikerscontrole, maar deze net versterkt.